\scanpage{8}\mspage{\handwritten{7}}
% Two lower corners of a drawn frame enclose the note above the first typed line.
\margin{\handwritten{Attention : la catégorie \uline{Hom} \unsure{gr.ab.|grab}, qui s'identifie à la catégorie définie par un $\unsure{\mathbb{R}|\mathrm{R}}\mathrm{Hom}_{\unsure{\mathbb Z|Z}}(C_{\unsure{*|\bullet}},C'_{\unsure{*|\bullet}})$ ; ce point devrait être élucidé !}}

de même les $(\mathrm{Cat})$-groupes(ou groupes à valeurs dans $(\mathrm{Cat})$) comme étant les groupes semisimpliciaux qui ``sont''des catégories i.e. qui \deleted{se définissent} satisfont la condition d'exactitude familière, permettant d'en déterminer tous les composants à l'aide de $C_{o\deleted{\unsure{o}}}$ et $C_1$. Appliquant le tapis de Dold, dans le cas commutatif, on trouve que les $(\mathrm{Cat})$-groupes abéliens forment une catégorie équivalente à celle des complexes de chaines de groupes abéliens qui sont de longueur 1:
\[
0\handwritten{\leftarrow}C_o\handwritten{\leftarrow}C_1\handwritten{\leftarrow}0\ .
\]

Quand on explicite la $(\mathrm{Cat})$-groupe abélien associé à un tel complexe de chaines, on trouve ``ma'' construction chérie de la ``catégorie de Picard'' associée au complexe: objets ceux de $C_o$, flêches de $x$ dans $y$ les $z$ dans \unsure{$C_1$} tels que $dz=y-x$. On trouve donc à isomorphisme ppès tous les \unsure{groupes} abéliens de $(\mathrm{Cat})$ ainsi. Un complément important est que toute catégorie de Picard \deleted{(resp. et commutative)} est équivalente (au sens techn\unsure{iqque|ique} qu'on ne précise pas ici) à \deleted{la catégorie} un vrai $(\mathrm{Cat})$-groupe \deleted{(resp. $(\mathrm{Cat})$-groupe commutatif)}.
% slanted
\margin{\handwritten{Quillen dixit, mais cela \unsure{contredit} les résultats \unsure{de ma|du} \unsure{lettre|livre} \unsure{à} \unsure{Verdier|Notations}, \unsure{la|voir} \unsure{même|sur|fin} \unsure{que|la|de la} \illegible{1} ? Mais attention, \unsure{même|vrai} si on part d'une cat. de Picard commut., elle ne se \unsure{réalise} que par un $(\mathrm{Cat})$-gr. qui \uuline{non} \uline{\unsure{comm.|commut.}}}}
\margin{\handwritten{Quand on simplifie mod ``\unsure{Grab|Gr.ab}'' équivalence cela correspond à travailler dans la catégorie \uline{dérivée} des complexes, \unsure{seuls} invariants \unsure{sont} bien évidemment $\pi_0,\pi_1$).}}

On notera qu'on a une \deleted{\unsure{interp}} classification beaucoup plus \unsure{restreinte} que pour les H-espaces (resp. les H-espaces homotopiquement comm.) mod équivalence d'homotopie, de façon précise entre les groupes resp. les groupes abéliens de la catégorie homotopique. La raison en est qu'on est beau\deleted{\unsure{x|c}}coup plus restrictif dans le contexte catégorique pour le sens de \unsure{l'associativité} (resp. et de la commutativité), puisqu'on ne demande pas seulement l'existence d'isomorphismes d'associ\unsure{aitivité|ativité}\deleted{\unsure{,|$\times$}} (qui suffirait pour avoir un ``groupe dans \inserted{la catégorie homotopique} $(\mathrm{Cat})'\handwritten{\unsure{\cong|\simeq}}(\mathrm{Ssimpl})'$, modulo l'existence d'unités et d'\unsure{inverses} mais qu'on suppose \uline{donné\deleted{e}s} ces isomorphismes, de façon à satisfaire à certaines compatibilités: c'est ce qui permet, me semble-t-il, de prouver \unsure{qu'à} équivalence près, on peut supposer qu'il s'agit de vrais groupes dans $(\mathrm{Cat})$, \deleted{resp. de vraie commutativité.}
% slanted
\margin{\handwritten{\deleted{\unsure{pourquoi|prouver}}\ \(\bigl(\)}}% opening bracket beside the final two lines

% archive sheet number: 7
